\subsection{背包}
背包问题在容量限制下选择物品。以下体积与容量均为正整数，第i件物品体积为$v_i$、价值为$w_i$，容量为V。容量上限版本令f[j]表示总体积不超过j的最大价值，允许不选，初值全0。
\subsubsection{01背包}
每件物品最多选一次，转移为$f_j=\max(f_j,f_{j-v_i}+w_i)$。容量倒序遍历，使右侧仍是处理当前物品之前的状态。
% snippet:knapsack01
\begin{lstlisting}
// f[0..V] initially zero; v, w use 1-based indices.
for (int i = 1; i <= n; i++)
    for (int j = V; j >= v[i]; j--)
        f[j] = std::max(f[j], f[j-v[i]] + w[i]);
\end{lstlisting}
答案f[V]，时间$O(nV)$，空间$O(V)$。若改为正序，一件体积2的物品更新f[2]后，又能由f[2]更新f[4]，等于重复选择同一件。
\subsubsection{完全背包}
每种物品可以无限选。容量正序遍历，让f[j-v[i]]包含本轮已经选过当前物品的方案。
% snippet:knapsack_complete
\begin{lstlisting}
for (int i = 1; i <= n; i++)
    for (int j = v[i]; j <= V; j++)
        f[j] = std::max(f[j], f[j-v[i]] + w[i]);
\end{lstlisting}
初值仍全0，答案f[V]。时间$O(nV)$，空间$O(V)$。
\subsubsection{多重背包}
每种物品至多有$c_i$件。把数量拆成1、2、4、……及剩余数量的若干包，每包按一个01物品处理，包的体积和价值都乘以包内件数。
% snippet:knapsack_multiple
\begin{lstlisting}
for (int i = 1; i <= n; i++) {
    int left = count[i];
    for (long long block = 1; left > 0; block *= 2) {
        int take = (int)std::min(block, (long long)left);
        left -= take;
        long long volume = 1LL * take * v[i];
        long long value = 1LL * take * w[i];
        if (volume > V) continue;
        for (int j = V; j >= volume; j--)
            f[j] = std::max(f[j], f[j-(int)volume] + value);
    }
}
\end{lstlisting}
例如13件拆成1、2、4、6件；前3包能组合成0到7件，再加6覆盖6到13件，因此0到13都能表示。count使用int，f使用long long；乘积及总价值须在long long范围内。时间$O(V\sum_i\log(c_i+1))$，空间$O(V)$。
\subsubsection{分组背包}
每组最多选一件。当前组的所有候选都从上一组的old数组转移，避免同组物品相互叠加。
% snippet:knapsack_group
\begin{lstlisting}
// group items are stored in v[g][k], w[g][k], k=0..size[g]-1.
for (int g = 0; g < groups; g++) {
    for (int j = 0; j <= V; j++) old[j] = f[j];
    for (int j = 0; j <= V; j++) {
        f[j] = old[j]; // Select none from this group.
        for (int k = 0; k < size[g]; k++)
            if (v[g][k] <= j)
                f[j] = std::max(f[j], old[j-v[g][k]] + w[g][k]);
    }
}
\end{lstlisting}
所有组共n件物品时，时间$O(nV)$，空间$O(V)$；若题目要求每组恰选一件，应删除“不选”的转移，并使用不可达初始化。
\subsubsection{恰好装满}
若f[j]表示总体积恰好j的最大价值，只能令$f_0=0$，其他容量初值为不可达。转移前检查来源是否可达。
% snippet:knapsack_exact
\begin{lstlisting}
const long long NEG = -(1LL << 60);
f[0] = 0;
for (int j = 1; j <= V; j++) f[j] = NEG;
for (int i = 1; i <= n; i++)
    for (int j = V; j >= v[i]; j--)
        if (f[j-v[i]] != NEG)
            f[j] = std::max(f[j], f[j-v[i]] + w[i]);
\end{lstlisting}
这里采用01背包，假定合法总价值绝对值小于$2^{60}$；f[V]仍为NEG表示无法恰好装满。不可达不是价值0，否则会凭空补足容量。
\subsubsection{背包计数}
计数状态ways[j]表示恰好组成体积j的方案数。初值ways[0]=1表示空选择，其他为0，转移使用加法。
% snippet:knapsack_count
\begin{lstlisting}
// Distinct positive coin types; MOD is a positive int.
ways[0] = 1;
for (int j = 1; j <= V; j++) ways[j] = 0;
for (int i = 1; i <= n; i++)
    for (int j = v[i]; j <= V; j++)
        ways[j] = (ways[j] + ways[j-v[i]]) % MOD;
\end{lstlisting}
ways使用long long。物品在外、容量正序，统计不限次数的组合，不区分取物品的先后顺序；若每件至多一次，容量改为倒序。容量在外、枚举最后一种物品在内，则统计有顺序的方案。面额1、2组成3时，组合为$\{1,1,1\}$与$\{1,2\}$共2种，有顺序则还区分$(1,2)$与$(2,1)$，共3种。时间$O(nV)$，空间$O(V)$。
